% Self-contained research fragment.  Requires amsmath, amssymb, amsthm.
% All labels use cr:.  References: BilinearIncoming, ChenSinger2012.
\section{Complex rational weights: a complete degree--rank theorem}
\label{cr:section}

The recent \emph{Bilinear Tails and Harmonic Determinants} contribution
proves that the formal harmonic degree of a moment determinant equals
its generic harmonic-pencil rank for every nonzero real rational weight.
It explicitly leaves the corresponding assertion for complex rational
weights open \cite[Section~5, Corollary \texttt{deg:rational};
Section~10, ``Complex weights and sharp cutoff exceptions'']{BilinearIncoming}.
We prove the complex assertion.  The new argument uses rationality in
the cutoff and the rigidity of symmetric rank-one updates.  It uses
ordinary transposition throughout; no Hermitian positivity is assumed.
The stronger numerical bound on the number of exceptions for eventually
one-sign real weights remains a separate advantage of the source theorem.

\subsection{A rational difference principle}

All fields in the next statements have characteristic zero.  The
indeterminate is $z$, and linear independence of functions means
independence over the constant field, unless another field is displayed.

\begin{lemma}[Equal ranks force equal images]\label{cr:rankone}
Let $A$ be a symmetric $q\times q$ matrix over a field $F$, let
$c\in F\setminus\{0\}$, and let $u\in F^q$.
If $u\notin\operatorname{im}A$, then
\[
 \operatorname{rank}(A+cuu^{\mathsf T})=\operatorname{rank}A+1.
\]
Consequently, if $A$ and $A+cuu^{\mathsf T}$ have equal rank, then
$u\in\operatorname{im}A$ and
$\operatorname{im}(A+cuu^{\mathsf T})=\operatorname{im}A$.
\end{lemma}
\begin{proof}
Put $d=\operatorname{rank}A$.  Symmetric elimination gives an
invertible matrix $S$ for which
\[
 S^{\mathsf T}AS=\begin{pmatrix}A_0&0\\0&0\end{pmatrix},
 \qquad \det A_0\ne0.
\]
Write $S^{\mathsf T}u=(v,w)^{\mathsf T}$.  If $u$ is outside
the image, some coordinate $w_j$ is nonzero.  The principal minor
on the $d$ coordinates of $A_0$ and this one further coordinate is
\begin{equation}\label{cr:minor}
 \det\begin{pmatrix}
 A_0+cvv^{\mathsf T}&cvw_j\\
 cw_jv^{\mathsf T}&cw_j^2
 \end{pmatrix}=cw_j^2\det A_0\ne0.
\end{equation}
For instance, subtract $(v_i/w_j)$ times the last row from row $i$
to verify the equality directly.  A rank-one update can increase
rank by at most one, proving the first claim.  Equal ranks therefore
force $u\in\operatorname{im}A$.  In that case the new image is
contained in the old one and has the same dimension, proving equality.
The proof remains valid when $w^{\mathsf T}w=0$; nonzero complex
isotropic vectors cause no exception to the argument.
\end{proof}

\begin{theorem}[Nonsingularity of rational symmetric primitives]
\label{cr:primitive}
Let $G(z)\in\operatorname{Sym}_q(F(z))$, let $w(z)\in F(z)$ be
nonzero, and let $p(z)\in F(z)^q$ have linearly independent
components over $F$.  If
\begin{equation}\label{cr:difference}
 G(z)-G(z-1)=w(z)p(z)p(z)^{\mathsf T},
\end{equation}
then $\det G(z)$ is a nonzero rational function.  This conclusion
holds for every symmetric rational solution, including after adding
an arbitrary constant symmetric matrix.
\end{theorem}
\begin{proof}
Suppose the generic rank were $d<q$.  A nonzero $d$-minor of $G$
has finitely many zeros and poles, and every $(d+1)$-minor vanishes
identically.  Hence there is an integer $N_0$ such that, at every
integer $N\ge N_0$, the matrices $G(N)$ and $G(N-1)$ are defined
and have rank $d$, while $w(N)\ne0$ and $p(N)$ is defined.
Characteristic zero ensures that these are infinitely many distinct
evaluation points.  Lemma~\ref{cr:rankone} applied to
\eqref{cr:difference} shows that the images of all these matrices
are one fixed $d$-dimensional subspace $W\subset F^q$, and that
$p(N)\in W$ at every sufficiently large integer.

Choose a nonzero linear functional $\lambda\in(F^q)^*$ vanishing
on $W$.  Then the rational function $\lambda p(z)$ vanishes at
infinitely many distinct points, so it is identically zero.  This
contradicts independence of the components of $p$.  Thus $d=q$.
Adding a constant symmetric matrix preserves \eqref{cr:difference},
so the same proof applies to every such addition.
\end{proof}

The result also applies with any fixed nonzero translation $h$ in
place of $1$, by changing variables.  Rationality is used to make
the ranks constant outside a finite set and to turn infinitely many
zeros into a function identity.  The statement is consequently stronger
than a positivity argument, but it does not make a claim about arbitrary
nonrational sequences of bilinear Gram matrices.

\subsection{The complex harmonic-degree theorem}

Let $C_1,\ldots,C_t\in\operatorname{Sym}_r(F)$ and put
\[
 C(X)=\sum_{j=1}^tX_jC_j,\qquad
 \rho=\operatorname{rank}_{F(X)}C(X),\qquad q=r-\rho.
\]
Suppose $R\in F(z)$ is nonzero and that the finite moment matrices
admit a decomposition
\begin{equation}\label{cr:decomposition}
 M_N=\sum_{n=0}^N R(n)v(n)v(n)^{\mathsf T}
     =D(N)+\sum_{j=1}^t h_j(N)C_j,
 \qquad v(n)=(1,n,\ldots,n^{r-1})^{\mathsf T},
\end{equation}
where $D(z)\in\operatorname{Sym}_r(F(z))$.  It suffices that the
decomposition and all its terms be defined for every sufficiently
large integer $N$.  The values $h_j(N)$ may lie in an extension of
$F$.  The variables $X_j$ are formal and independent; no arithmetic
independence of the actual harmonic values is asserted.

\begin{theorem}[Complex rational degree--rank equality]
\label{cr:degree}
Under these hypotheses,
\begin{equation}\label{cr:main}
 \boxed{\deg_X\det(D(z)+C(X))=\rho\quad\text{over }F(z).}
\end{equation}
Only finitely many admissible integer cutoffs can have smaller
degree.  In particular the conclusion holds for every nonzero complex
rational weight with no poles at nonnegative integers and its usual
finite harmonic decomposition.
\end{theorem}
\begin{proof}
The determinant has degree at most $\rho$, since any term choosing
more than $\rho$ columns from $C(X)$ is zero.  Since $F$ is infinite,
choose $a\in F^t$ for which $A=C(a)$ has rank $\rho$.  Let $K$
be an $r\times q$ constant matrix whose columns form its kernel.
Choose a complement $U$ so that $S=[U\ K]$ is invertible and
\[
 S^{\mathsf T}AS=\begin{pmatrix}A_0&0\\0&0\end{pmatrix},
 \qquad\det A_0\ne0.
\]
Every matrix $A+tC_j$ has rank at most $\rho$.  Over $F(t)$
its upper-left block is invertible near $t=0$, so its Schur complement
vanishes identically.  The coefficient of $t$ in that complement is
$K^{\mathsf T}C_jK$.  Therefore
\begin{equation}\label{cr:isotropy}
 K^{\mathsf T}C_jK=0\qquad(1\le j\le t).
\end{equation}
This algebraic argument is the complex form of the source's
compression lemma.

Put $G(z)=K^{\mathsf T}D(z)K$ and $p(z)=K^{\mathsf T}v(z)$.
The components of $p$ are independent polynomials, because $K$
has column rank $q$.  Equations \eqref{cr:decomposition} and
\eqref{cr:isotropy} show, at all sufficiently large integers, that
\[
 G(N)-G(N-1)=R(N)p(N)p(N)^{\mathsf T}.
\]
Both sides are rational functions in $N$, so this is the identity
\eqref{cr:difference}.  Theorem~\ref{cr:primitive} gives
$\det G(z)\ne0$.  Finally,
\begin{equation}\label{cr:direction}
 [t^\rho]\det(D(z)+tA)
   =\frac{\det A_0}{(\det S)^2}\det G(z)\ne0,
\end{equation}
obtained by selecting $tA_0$ in its $\rho$ rows and columns after
congruence by $S$.  This proves \eqref{cr:main}.  The same
coefficient is nonzero at all but finitely many admissible integer
cutoffs, proving the final assertion.  When $q=0$, its determinant
is one and the conclusion is immediate.
\end{proof}

No common kernel of the entire pencil is assumed.  The kernel of
one maximal-rank member is only totally isotropic for every $C_j$,
which is exactly the compression property needed in the proof.
The theorem answers the complex-weight question affirmatively rather
than restricting it to pencils with a common kernel.

\subsection{Finite certificates for every exceptional cutoff}

The proof gives a cheap sufficient certificate using just one
maximal-rank direction.  Choose a nonzero polynomial $d(z)$ clearing
the denominators of $G(z)$ and set $H(z)=d(z)G(z)$.  Then
\begin{equation}\label{cr:certificate}
 0\ne E(z):=\det H(z)\in F[z].
\end{equation}
Every degree-deficient cutoff at which $d(N)\ne0$ is a root of
$E$.  If $e_{ij}=\deg H_{ij}$, with degree $-\infty$ for zero
entries, then
\begin{equation}\label{cr:degree-bound}
 \deg E\le \max_{\pi\in S_q}\sum_{i=1}^q e_{i,\pi(i)}
 \le q\max_{i,j}e_{ij}.
\end{equation}
Thus the right side bounds the number of distinct such exceptions.
This is a bound computable from rational entries and does not depend
on ordering the coefficient field.  Apparent poles introduced by
denominator clearing must be removed or checked in the original
decomposition.

For an exact rather than sufficient certificate, write the highest
homogeneous component after clearing a common denominator in $z$ as
\[
 [\det(D(z)+C(X))]_\rho
       =\frac{1}{b(z)}\sum_{|\nu|=\rho}e_\nu(z)X^\nu.
\]
Theorem~\ref{cr:degree} guarantees that not all $e_\nu$ vanish.
For an admissible $N$ with $b(N)\ne0$, degree is deficient exactly
when $g(N)=0$, where
\begin{equation}\label{cr:gcd}
 g(z)=\gcd_{|\nu|=\rho}e_\nu(z).
\end{equation}
This is the complex-field extension
of the source's coefficient-gcd test; it is not a new polynomial-gcd
algorithm.  The cheaper polynomial $E$ supplies a finite candidate
set, and the full highest component filters it exactly.

\subsection{A residue interpretation of the rational compression}

For clarity, the usual harmonic decomposition is available over
$\mathbb C$ without a sign assumption.  Partial fractions of
$z^mR(z)$ consist of a polynomial and finitely many terms
$c/(z-\alpha)^j$.  Poles differing by integers belong to one shift
orbit.  Summing those terms from $0$ to $N$ produces one generalized
harmonic coordinate per orbit and pole order, together with rational
endpoint corrections.  This is the form \eqref{cr:decomposition}.

The discrete residue of a rational function at an orbit and pole
order is the sum of the corresponding partial-fraction coefficients
over that orbit.  The classical rational-summability criterion states
that a rational function has a rational antidifference if and only if
all these discrete residues vanish \cite[Proposition~2.5]{ChenSinger2012}.
For a polynomial subspace $V$ with basis $p_1,\ldots,p_q$, the
following statements are equivalent:
\begin{enumerate}
\item every $R(z)p_i(z)p_j(z)$ has a rational antidifference;
\item every discrete-residue bilinear form
$(f,g)\mapsto\operatorname{dres}_{\mathcal O,k}(Rfg)$ vanishes
on $V\times V$;
\item if $K$ is the coefficient matrix of this basis, then
$K^{\mathsf T}C_jK=0$ for every harmonic coefficient matrix.
\end{enumerate}
The equivalence is the classical residue criterion applied entry by
entry.  Theorem~\ref{cr:primitive} adds the following consequence:
whenever these conditions hold and the $p_i$ are independent, every
symmetric rational matrix primitive of the products $Rp_ip_j$ is
nonsingular over $\mathbb C(z)$.  Thus a maximal-rank pencil kernel
produces a rationally summable polynomial space whose accumulated
bilinear form cannot be identically degenerate.

\subsection{Exact nonreal examples}

Let
\[
 R_\lambda(z)=\frac1{z+1}+\frac\lambda{z+2},\qquad\lambda\ne0,
 \qquad X=\sum_{n=0}^N\frac1{n+1}.
\]
For an $r\times r$ moment matrix, $r\ge2$, the only harmonic
coefficient matrix is
\[
 C=u_1u_1^{\mathsf T}+\lambda u_2u_2^{\mathsf T},\qquad
 u_a=(1,-a,\ldots,(-a)^{r-1})^{\mathsf T},
\]
and has rank two.  Thus the harmonic degree is exactly two over
$\mathbb C(N)$, for every nonzero complex $\lambda$.

At $r=3$, the kernel polynomial is $(z+1)(z+2)$ and
\begin{equation}\label{cr:three}
 \boxed{
 [X^2]\det(D(N)+XC)=
 \frac{\lambda}{12}(N+1)(N+2)(N+3)
       \{3(\lambda+1)N+4\lambda+8\}.}
\end{equation}
At $r=4$, use the two kernel polynomials
$(z+1)(z+2)$ and $z(z+1)(z+2)$.  Then
\begin{equation}\label{cr:four}
 \begin{split}
 [X^2]\det(D(N)+XC)
 ={}&\frac{\lambda}{3600}N(N+1)^2(N+2)^2(N+3)^2(N+4)\\
 &\quad\cdot\bigl\{6(\lambda+1)^2N^2
 +(14\lambda^2+48\lambda+34)N
 +10\lambda^2+45\lambda+50\bigr\}.
 \end{split}
\end{equation}
Both identities follow from \eqref{cr:direction} and finite sums of
polynomials.  In the four-dimensional case, the two-by-two compressed
determinant is the expression after the initial factor $\lambda$.
For $\lambda=i$, its final quadratic is
$20N+40+i(12N^2+48N+45)$, so it cannot vanish at a nonnegative
integer.  Consequently the four-dimensional harmonic degree is
exactly two at every $N\ge1$, even though the actual evaluated
moment determinant vanishes at $N=1,2$ because too few nodes are
present.  Formal degree and numerical determinant value remain
different notions.

A nonreal rational weight can nevertheless have a genuine later
degree exception.  Take $r=4$ and
\begin{equation}\label{cr:exception-weight}
 R(z)=\frac{39}{z+1}+\frac{78i}{z+2}
                        -\frac{61+95i}{z+3}.
\end{equation}
Its one harmonic matrix has rank three.  The kernel polynomial is
$p(z)=(z+1)(z+2)(z+3)$, and exact Faulhaber summation gives
\begin{equation}\label{cr:exception-compression}
 \begin{split}
 \sum_{n=0}^N R(n)p(n)^2
  =-\frac16&(N-3)(N+1)(N+2)(N+3)(N+4)\\
           &\cdot\{(22+17i)N+56+22i\}.
 \end{split}
\end{equation}
The top-left three-square minor of the harmonic matrix is
$1155960-742248i\ne0$.  Hence the only degree-deficient
nonnegative cutoff is $N=3$.  At that cutoff the formal determinant
is exactly
\begin{equation}\label{cr:exception-polynomial}
 \begin{split}
 P_3(X)={}&(-22046418576+2049048144i)X^2\\
 &+\left(\frac{455872142232}{5}-7792299840i\right)X\\
 &-\frac{2356478577592}{25}
                         +\frac{183512850456}{25}i.
 \end{split}
\end{equation}
Thus its degree drops from three to two, while substitution of the
actual $X=25/12$ gives
\[
 P_3(25/12)=\frac{55033-146419i}{25}\ne0.
\]
Independently, this value is the four-node Vandermonde square times
$R(0)R(1)R(2)R(3)$.  The exception is therefore a formal-degree
phenomenon at a nonzero evaluated determinant, and not a lack of
nodes or a vanishing weight.

\subsection{Integration and further questions}

The source's real-weight results are valid and need no retraction.
Its open complex-weight question can be marked solved under the
same rational-decomposition hypotheses, and Corollary
\texttt{deg:rational} can be broadened to complex rational weights.
The eventual-sign bound $r(r-\rho)$ should retain its original real
hypotheses.  The present denominator-degree certificate is a
different bound and is not asserted to equal that real bound.

Two precise next questions remain.  First, determine sharp bounds
for the exceptional-cutoff polynomial using only the pole orders and
integer dispersions of $R$, without constructing the whole compressed
primitive.  Second, classify polynomial subspaces totally isotropic
for all discrete-residue forms of $Rfg$; the primitive theorem proves
nonsingularity once such a subspace is given, but does not classify
the subspaces or the rank-two pencils producing them.

